% Tipo de documento. En este caso es un art�culo, para folios A4, tama�o de la fuente 11pt y con p�gina separada para el t�tulo
\documentclass[a4paper,11pt,titlepage]{article}

% Carga de paquetes necesarios. OrdenesArticle es un paquete personalizado
\usepackage[spanish]{babel} 
\RequirePackage[T1]{fontenc}
\RequirePackage[ansinew]{inputenx} 
\usepackage[spanish,cap,cont,title,fancy]{OrdenesArticle}
\usepackage{lmodern}
\usepackage{array}
\usepackage{graphicx}
\usepackage{hyperref}
\usepackage{pifont}
\usepackage{listings}
\usepackage[usenames,dvipsnames]{color}
\usepackage{colortbl}
\usepackage{color}
\usepackage{ifthen}
\usepackage{longtable}
\hypersetup{bookmarksopen,bookmarksopenlevel=4,linktocpage,colorlinks,urlcolor=blue,citecolor=blue,
						linkcolor=blue,filecolor=blue,pdfnewwindow,
						pdftitle={PlanCoDE: Planificador Colaborativo para el Dise�o de Estrategias y Acciones de Emergencia},
						pdfauthor={Juan Andrada Romero, Jose Domingo L�pez L�pez},
						pdfsubject={Sistemas para la Colaboraci�n}}


% Macro para definir una lista personalizada 
\newenvironment{milista}%
{\begin{list}{\textbullet}%
{\settowidth{\labelwidth}{\textbullet} \setlength{\leftmargin}{\dimexpr\labelsep+\labelwidth+5pt}
\setlength{\itemsep}{\dimexpr 0.5ex plus 0.25ex minus 0.25ex}
\setlength{\parsep}{\itemsep}
\setlength{\partopsep}{\itemsep}
\addtolength{\topsep}{-7.5pt}
}}%
{\end{list}}

% Macro para insertar una imagen
%       Uso: \imagen{nombreFichero}{Factor escala}{Caption (leyenda)}{Label (identificador para referenciarla)}
% -------------------------------------------------------------------------------------------------------------
\def\imagen#1#2#3#4{
 \begin{figure}[h]
 \begin{center}
   \scalebox{#2}{\includegraphics{#1}}
 \caption {#3}
 \label{#4}
 \end{center}
 \end{figure}
}

% Macro para dar color a las filas de la tabla "`Operaciones"'
\def\colorFila{
	\rowcolor[rgb]{0.4,0.52,1.0}
}

% Definicion del "listings" para el lenguaje Java
\lstdefinelanguage{Java}
{
 morecomment = [l]{//}, 
 morecomment = [l]{///},
 morecomment = [s]{/*}{*/},
 morestring = [b]", 
 sensitive = true,
 morekeywords = {package, static, while, switch, break, line, void, String, Object, int, Integer, instanceof, else, if, for, private, return, new, public, class, import, int, boolean, true, false, extends, final, super, protected, abstract, this, do, float, double, null, try, catch, implements}
}

% Configuraci�n del c�digo incluido para el lenguaje Java
\lstset{
  language=Java,
  basicstyle=\footnotesize,
  backgroundcolor=\color{white},
  showspaces=false,
  showstringspaces=false,
  showtabs=false,
  frame=single,
  tabsize=2,
  captionpos=b,
  breaklines=true,
  breakatwhitespace=false,
  escapeinside={\%},
  keywordstyle = \color [rgb]{0,0,1},
  commentstyle = \color [rgb]{0.133,0.545,0.133},
  stringstyle = \color [rgb]{0.627,0.126,0.941}
}

\begin{document}

% En las p�ginas de portada e �ndices, no hay encabezado ni pie de p�gina
\pagestyle{empty} 

% Se incluye la portada
\begin{titlepage}
	\begin{center}
  	{\LARGE UNIVERSIDAD DE CASTILLA-LA MANCHA} \\
  	\bigskip
  	{\Large ESCUELA SUPERIOR DE INFORM�TICA} \\
  	\vspace{20mm}
  	\includegraphics[scale=0.45, keepaspectratio]{./images/esi_bw.png} \\
  	\vspace{20mm}
  	{\Huge \textbf{Almacenamiento y Recuperaci�n de la Informaci�n}}\\
  	\vspace{10mm}
  	{\LARGE \textsc{\textbf{M�dulo de importaci�n e indexaci�n de documentos}}}\\
  	\vspace{30mm}
  	\today
  	\vspace{40mm}
  	\begin{flushleft}
  		{\large Juan Andrada Romero}\\
  		\vspace{1mm}
  		{\large Jose Domingo L�pez L�pez}\\		
  	\end{flushleft}
	\end{center}
\end{titlepage}
\clearpage

% Texto del reverso de la portada
%%\mbox{}
%%\vspace{18cm}
%%\begin{small}
% Se ajusta la separaci�n entre p�rrafos
%%\parskip=10pt 

%%\copyright~ 2008/2009 Juan Andrada Romero. Universidad de Castilla La Mancha, Escuela Superior de Inform�tica de Ciudad Real.

%%Se permite la modificaci�n, copia y distribuci�n de este documento, seg�n la licencia de documentaci�n GNU (\url{http://www.gnu.org}).

%%Este documento fue compuesto con \LaTeX{}. Im�genes generadas con Power Point y Gimp.
%%\end{small}

%%\newpage

% En las p�ginas de �ndices y prefacio, se utiliza numeraci�n romana
\pagenumbering{Roman}

% Se crea el �ndice
\tableofcontents
% Se pasa p�gina y se a�ade el �ndice de figuras al �ndice principal
%\clearpage\phantomsection
%\addcontentsline{toc}{section}{\listfigurename}
% Se crea el �ndice de figuras
%\listoffigures
% Si se quiere crear un �ndice de tablas se pondr�a: \listoftables

\newpage

% Se ajusta la separaci�n entre p�rrafos
\parskip=10pt

% Se a�ade el prefacio al �ndice
%\clearpage\phantomsection
%\addcontentsline{toc}{section}{Prefacio}

% Comienza el contenido del documento. Se utilizan n�meros ar�bigos y el encabezado y pie de p�gina personalizado
\pagenumbering{arabic}
\pagestyle{fancy}

\section{Introducci�n}


\section{Interfaz de usuario}
La interfaz gr�fica de usuario de la herramienta PlanCoDE es de tipo WYSIWIS (\textit{What You See Is What I see}) relajado. La raz�n por la que se ha elegido la variante \textit{relajada} es para evitar los problemas asociados con las interfaces WYSIWIS estrictas, tales como las guerras de ventanas y la sobrecarga de informaci�n. A continuaci�n se profundizar� sobre algunas de las caracter�sticas que se han relajado.

\subsection{Separaci�n del espacio de trabajo} \label{sec:separacion-workspace}
Se ha relajado la barra de men�s y las ventanas de uso privado, tales como las ventanas de configuraci�n o carga de planos/im�genes, as� como la caja de herramientas. En definitiva, las partes sobre las que se aplica el WYSIWIS es sobre el log del chat, el log de las acciones, el �rea de dibujo y la lista de usuarios activos en la sesi�n. De este modo, los usuarios pueden colaborar para realizar el dise�o de las estrategias que estimen oportunas sin perder la intimidad y la privacidad de uso de la aplicaci�n.

\subsection{Visualizaci�n de cursores}
Cada usuario tiene su propio cursor privado y ning�n otro usuario puede observarlo ni manipularlo. Esto es debido a que PlanCoDE ser� portada, en un futuro, a dispositivos m�viles. Este tipo de pantallas son muy reducidas y no es recomendable permitir que los cursores de los distintos usuarios las invadan, causando confusi�n. No obstante, hay un caso especial de puntero que s� ser� visualizado por todos los usuarios (ver secci�n \ref{sec:telepunteros}).

\subsection{Gesti�n de la distribuci�n en la pantalla}
La interfaz gr�fica de usuario de la aplicaci�n de escritorio que se har� como primera entrega es un prototipo, por lo que esta caracter�stica no se satisfar� por completo. No obstante, como se puede deducir de la secci�n \ref{sec:separacion-workspace}, cada usuario podr� elegir d�nde situar las ventanas \textit{no compartidas} de la aplicaci�n de forma privada.

\subsection{Gesti�n de la informaci�n visualizada}
La aplicaci�n dispone de una ventana de log donde se ir�n registrando todas las acciones que realizan cada uno de los distintos usuarios en el �rea de dibujo. De este modo se representa de forma textual la informaci�n que se va generando de forma gr�fica.

\subsection{Acoplamiento de interfaces de usuario}
El nivel de acoplamiento (ver Figura \ref{fig:acoplamiento}) entre las interfaces de usuario de los distintos clientes se clasifica como \textit{medio}. Esto es debido a que hay partes de la interfaz gr�fica que son privadas cada usuario, pero hay otras partes como el chat, la lista de usuarios activos y el �rea de dibujo que son compartidos (ver secci�n \ref{sec:separacion-workspace}).

\imagen{imagenes/acoplamiento.png}{0.4}{Niveles de Acoplamiento de Interfaces de Usuario}{fig:acoplamiento}

\subsection{Telepunteros} \label{sec:telepunteros}
Una caracter�stica muy importante de la aplicaci�n es el uso de un telepuntero. Un telepuntero se utiliza para apuntar informaci�n compartida en ventanas p�blicas y s�lo una persona puede tener acceso al mismo (acceso mutuamente excluyente).

Todos los usuarios ver�n en su pantalla dos cursores: el privado y el telepuntero. Los usuarios de PlanCoDE pueden utilizar el telepuntero, previamente solicitado y concedido, para realizar acciones sobre el �rea de dibujo (dibujar, borrar o colocar elementos). Este mecanismo se implementa para evitar que todos los usuarios traten de dibujar o borrar al mismo tiempo provocando confusi�n al resto de usuarios.

Adem�s, al lado del telepuntero se mostrar� el nombre del usuario que lo tiene en posesi�n y la acci�n que va a realizar sobre el �rea de dibujo. De este modo, los usuarios que observan el telepuntero pueden tener una percepci�n clara de \textit{qui�n} est� realizando \textit{qu�} en cada momento.
\section{Soporte de \textit{awareness}}

En esta secci�n se va comentar el soporte \textit{awareness} del espacio de trabajo de la aplicaci�n, es decir, como se soporta la interacci�n entre los diferentes usuarios de la aplicaci�n. 
As�, se van detallar las t�cnicas empleadas para intentar responder a las preguntas que aparecen en \cite{tema5}. Cabe destacar que que no es necesario responder a todas ellas, sobrecargando incluso el sistema si se responden.

\begin{milista}
	\item \textbf{Identidad}: PlanCoDE mantendr� una lista de usuarios conectados, indicando su rol dentro del sistema. As�, se indicar� si el usuario corresponde a un cuerpo policial, a un cuerpo de bomberos, etc. De este modo, el resto de usuarios ser�n conscientes de qu� tipos de usuarios est�n haciendo uso de la herramienta. 
	Adem�s, cada usuario, y todas las acciones que �ste realice, se van a identificar por un color.
	\item \textbf{Acciones}: para indicar lo que cada usuario est� haciendo, se va a utilizar el concepto de \textit{telepuntero}, como ya se ha comentado en la secci�n \ref{sec:telepunteros}. 
	\indent Por otra parte, cuando un usuario comience a escribir en el chat, el resto de usuarios tambi�n ser�n informados de esta situaci�n, ya que se mostrar� un mensaje indicando que ese usuario est� escribiendo texto.
	\item \textbf{Intenciones}: cuando un usuario quiera utilizar el �rea de dibujo, debe solicitar el \textit{telepuntero}. De este modo, se informar� al resto de usuarios acerca de qui�n lo ha solicitado, pudiendo saber qui�n va a manipular el �rea de trabajo a continuaci�n.
	\item \textbf{Cambios}: relacionado con todo lo anterior, un usuario puede conocer los cambios que otro usuario est� haciendo sobre el �rea de dibujo, pues dichos cambios se ir�n mostrando al utilizar el \textit{telepuntero}. \\
	\indent Adem�s, todos los cambios realizados en el �rea de trabajo se almacenar�n tambi�n en forma de \textit{log} textual, por lo que cualquier usuario puede consultar las acciones que ha llevado a cabo otro usuario.
	\item \textbf{Objetos}: el objeto que se utiliza por parte de los usuarios es el \textit{telepuntero}, ya que para poder dibujar en el �rea de trabajo, previamente deben solicitar su turno para poder hacerlo.
\end{milista}

Para terminar, cabe destacar que para dotar a PlanCoDE de mayor funcionalidad, se ha incluido un canal adicional de texto, que es el chat, para que los usuarios puedan comunicarse en tiempo real. Es un canal adicional porque el principal objetivo de esta herramienta es el dise�o de un plan de emergencia realizando modificaciones sobre un �rea de dibujo, en la cu�l se carga una imagen o un mapa.

% Secci�n de ap�ndices
% Se a�ade el titulo de "Ap�ndices" al �ndice
%\addcontentsline{toc}{section}{Ap�ndices}

%\input{Apendices.tex}

% A�adimos la bibliografia al �ndice
\clearpage\phantomsection
\addcontentsline{toc}{section}{Referencias}

% Bibliograf�a. En este caso se usa BibTeX
\bibliographystyle{plain}
\pagestyle{plain} 
\bibliography{Bibliografia}

\end{document}
